\documentclass[11pt]{article}
\usepackage{../../latex/stat4ml_style}

\title{Chapter 1: Learning Algorithms Overview \\ \large Part II --- Introduction to Statistical Learning}
\author{Omid Safarzadeh}
\date{}

\begin{document}
\maketitle

\begin{abstract}
\noindent
Part I gave us the language of probability -- sample spaces, random
variables, moments. Part II turns that language into \emph{learning
algorithms}: procedures that look at data and produce a function, a
partition, or a policy. This opening chapter is a map of the territory: the
handful of learning paradigms (supervised, unsupervised, semi-supervised,
online, reinforcement, graph-based, and distillation-based) that every
later chapter in this part is an instance of. Read it as a table of
contents with explanations, not as a self-contained technical chapter --
each paradigm gets its full mathematical treatment later (regression and
cross-validation in Chapter 2, unsupervised learning and PCA in Chapter 5,
and so on).

\emph{Acknowledgement: this chapter expands the Introduction to Statistical
Learning lecture slides by Omid Safarzadeh. Course texts: James, Witten,
Hastie \& Tibshirani, \emph{An Introduction to Statistical Learning}; Murphy,
\emph{Machine Learning: A Probabilistic Perspective}.}
\end{abstract}

\tableofcontents

\section{Supervised Learning}

\begin{definition}
In supervised learning we observe \emph{labeled} training data
$\{(x_i, y_i)\}_{i=1}^n$, where $x_i$ is an instance and $y_i$ its label, and
we assume a generating model
\[
  y = f(x) + \epsilon,
\]
where $\epsilon$ is unobserved noise. A learning algorithm maps the training
set to an estimate $\hat f(\cdot)$ of $f$. The goal is that $\hat f(x)$ is
close to $y$ for \emph{new} $(x,y)$ pairs, not just the ones we trained on --
this distinction (fit vs.\ generalize) is the entire subject of Chapter 2
(bias--variance tradeoff, cross-validation).
\end{definition}

\begin{example}[Advertising]
$X$ = advertising budget (\$000s), $Y$ = sales (000s of units); each data
point is one past advertising campaign. Fitting a line to this scatterplot
by least squares is the simplest possible supervised learning algorithm --
formalized fully in Chapter 2.
\end{example}

\begin{example}[Image classification -- CIFAR-10]
$X$ = a $32\times32$ color image, $Y \in \{0,\dots,9\}$ one of ten object
classes. The same supervised-learning definition applies; only $f$ (now a
convolutional network, Part II Chapter 9) and the loss function (cross-
entropy, not squared error) change. (Krizhevsky \& Hinton, 2009.)
\end{example}

\section{Unsupervised Learning}

\begin{definition}
In unsupervised learning we observe \emph{unlabeled} data $x_1,\dots,x_n$
(no $y_i$) and ask what structure -- clusters, a lower-dimensional
representation, a density -- the data itself reveals.
\end{definition}

\begin{example}[Clustering]
Given unlabeled points, $K$-means with $K=3$ partitions them into three
groups by proximity to three learned centroids (full treatment: Part II,
Chapter 7).
\end{example}

\begin{example}[Learned image compression]
Image compression represents an image with as few bits as possible.
Learned, probability-model-based codecs have been shown to beat classical
codecs (BPG, JPEG, JPEG2000) on perceptual quality at the same bit rate
(Mentzer et al., 2018) -- an unsupervised (or self-supervised) learning
problem: there is no external label, only the image itself to reconstruct.
\end{example}

\begin{example}[Masked language modeling]
BERT (Devlin et al., 2018) is trained by masking tokens in a sentence and
predicting them from context -- the ``label'' is manufactured from the
unlabeled text itself. This blurs the line between unsupervised and
supervised learning, which is exactly why the next paradigm exists.
\end{example}

\section{Semi-Supervised Learning}

\begin{definition}
Semi-supervised learning uses a small labeled dataset together with a much
larger unlabeled dataset, typically by having a model trained on the
labeled data generate ``pseudo-labels'' for the unlabeled data, which are
then used for further training.
\end{definition}

\begin{example}
Billion-scale semi-supervised image classification: a teacher model labels
a huge unlabeled image corpus; a student model trains on the combined
labeled + pseudo-labeled data (Yalniz et al., 2019).
\end{example}

\subsection{Further Examples Worth Knowing}
\begin{itemize}[nosep]
  \item \textbf{Deep belief networks} for handwritten-digit classification
        and generation -- an early (pre-2012) example of layerwise
        unsupervised pretraining feeding into supervised fine-tuning.
  \item \textbf{Reinforcement learning} milestones: emergent tool-use in
        multi-agent hide-and-seek, and AlphaGo / AlphaZero learning to play
        Go (and chess/shogi) purely from self-play.
\end{itemize}

\section{Online Learning}

Ordinary supervised/unsupervised learning assumes: data arrives as one
static batch; the model is trained once; the model is then frozen. Three
problems follow directly from that assumption:
\begin{itemize}[nosep]
  \item it is not well-suited to genuinely streaming data;
  \item retraining from scratch on the full history is wasteful once the
        history is large;
  \item the fitted model becomes stale as the data-generating process
        drifts over time (\emph{concept drift}).
\end{itemize}

\begin{definition}
\emph{Online (incremental) learning} algorithms update their parameters one
new observation (or mini-batch) at a time, without revisiting the full
history, so the model can track a data stream as it evolves.
\end{definition}

The simplest possible instance is an incremental (running) mean: given the
current mean $\bar x_{n-1}$ over $n-1$ points and a new observation $x_n$,
\[
  \bar x_n = \bar x_{n-1} + \frac{x_n - \bar x_{n-1}}{n},
\]
which is mathematically identical to the batch mean but needs $O(1)$ memory
and work per new point instead of $O(n)$. Exercise 1 asks you to implement
exactly this update (Putatunda, 2021).

\section{Reinforcement Learning}

\begin{definition}
In reinforcement learning an \emph{agent} interacts with an
\emph{environment} over a sequence of time steps, choosing actions to
maximize cumulative reward, with no explicit $(x_i, y_i)$ pairs telling it
the ``correct'' action -- only a reward signal, often delayed. This is
covered here only as a paradigm; a full treatment (Markov decision
processes, policy gradients, multi-armed bandits) is planned for a future
part of this book.
\end{definition}

\section{Graph Representation Learning}

\begin{definition}
When data is naturally a graph (a social network, a molecule, a knowledge
graph), graph representation learning produces vector embeddings of nodes
(or whole graphs) that preserve graph structure, so that ordinary
supervised or unsupervised algorithms can be applied downstream. This
paradigm is also deferred to a full treatment in Part III.
\end{definition}

\section{Knowledge Distillation}

\begin{definition}
In \emph{knowledge distillation} (teacher/student learning), a large
``teacher'' model's output distribution (not just its hard labels) is used
as a training target for a smaller ``student'' model, transferring the
teacher's learned function into a cheaper model to deploy. This is a
supervised-learning technique whose ``labels'' are themselves the output of
another learned model -- a nice illustration of how the paradigms in this
chapter combine rather than partition the field.
\end{definition}

\section{Summary}
\begin{summarybox}
\begin{itemize}[nosep]
  \item \textbf{Supervised}: learn $\hat f$ from labeled pairs $(x_i, y_i)$.
  \item \textbf{Unsupervised}: find structure in unlabeled $x_i$ alone.
  \item \textbf{Semi-supervised}: combine a little labeled data with a lot
        of unlabeled data via pseudo-labeling.
  \item \textbf{Online}: update the model incrementally as data streams in,
        to track concept drift.
  \item \textbf{Reinforcement}: learn a policy from reward signals, not
        labels.
  \item \textbf{Graph representation learning}: embed graph-structured data
        for downstream use.
  \item \textbf{Knowledge distillation}: train a small model to imitate a
        large one's output distribution.
\end{itemize}
These are not mutually exclusive categories -- masked language modeling
(BERT) and knowledge distillation both blend supervised mechanics onto
unsupervised or model-generated targets. The rest of Part II works through
each paradigm's core algorithms in full mathematical detail.
\end{summarybox}

\section{Exercises}
See \texttt{exercises/python/exercise01.py}, \texttt{exercises/r/exercise01.R},
and \texttt{exercises/rust/exercise01.rs}: a closed-form ordinary-least-
-squares fit (supervised), a single $K$-means assignment step
(unsupervised), and the incremental-mean update (online learning) --
one small, working example of each of the three paradigms most fully
developed later in Part II. The broken counterparts to debug are in
\texttt{exercise01\_buggy.\{py,R,rs\}}.

\section*{References}
\begin{itemize}[nosep]
  \item Devlin, J., Chang, M.-W., Lee, K., \& Toutanova, K. (2018). BERT:
        Pre-training of deep bidirectional transformers for language
        understanding. \emph{arXiv:1810.04805}.
  \item Krizhevsky, A., \& Hinton, G. (2009). Learning multiple layers of
        features from tiny images.
  \item Mentzer, F., Agustsson, E., Tschannen, M., Timofte, R., \& Van Gool,
        L. (2018). Conditional probability models for deep image
        compression. \emph{CVPR}.
  \item Putatunda, S. (2021). \emph{Practical Machine Learning for Streaming
        Data with Python}.
  \item Yalniz, I.~Z., J\'egou, H., Chen, K., Paluri, M., \& Mahajan, D.
        (2019). Billion-scale semi-supervised learning for image
        classification. \emph{arXiv:1905.00546}.
\end{itemize}

\end{document}
